\section{Related Work}

Four lines of work touch the question this paper asks, and each answers a different part of it. Equivariance measurement scores how much of a transformation a representation carries; symmetry-based architectures build that transformation in; colour and scale-space theory fix what the reference transformations analytically are; and causal abstraction asks when an intervention on a representation corresponds to one on the world. No single line contains the join this paper makes --- existence conditions, realisation form, composition price, and construction --- and \S{}2.5 states the gap line by line.

\subsection{Transformations in existing representations}

The oldest question in this line is whether a representation has \textit{learned} the transformation it is meant to respect. Lenc and Vedaldi [1] measure equivariance and equivalence by fitting a probe-space linear map and scoring how well it commutes with the transformation, which is also the object this paper fits as a baseline operator. The line has since become quantitative: the Lie derivative [2] scores \textit{local} equivariance error along a transformation's flow; Bruintjes et al. [3] ask what governs learned equivariance across layers and find architecture-level dependence; Romero and Lohit [4] learn partial equivariances rather than assuming a group; Brehmer et al. [5] ask whether equivariance matters at scale; and He et al. [6] trace how data augmentation shapes the representations themselves. Equivariant self-supervision [7], [8] identifies or encourages equivariant embeddings rather than measuring them after the fact.

What this line supplies is the empirical object: a frozen representation whose response to a known transformation can be measured, and the practice of grading that response against a re-rendered reference. What it does not contain is an \textit{information condition}. Equivariance is scored, never asked to exist; the score conflates the part of the algebra the representation keeps with the part it merges away; and no score is defined so that its composition over a chain is meaningful. Our fibre condition supplies the missing information condition --- an induced action exists exactly when the encoder's fibres survive the transformation, and existence then delivers the composition law for free --- and \S{}5.4 grades a chain against the truth-level composite rather than treating it as a test of a group axiom. The measurement-side consequence is sharp: because the scores presuppose an operator family, a family can be self-consistent and still fail the transformation, and the standard proxies can rank an unfaithful operator first (\S{}5.5).

\subsection{Symmetry-based and equivariant representations}

A second line builds the symmetry in. Group-equivariant convolutions [9], steerable variants [10], and their colour-specific descendants --- CEConv [11], learning colour-equivariant representations [12], and exact colour equivariance by hypertoroidal covering [13] --- construct layers whose response to the group is exact by design. A parallel line studies what such networks can express and when pointwise activations preserve the symmetry [14], and symmetry-based disentanglement defines group-structured latent factors from the data [15]. Classical representation theory [16] supplies the harmonic-block decomposition this paper uses as its candidate \textit{forms} (\S{}3.6).

These methods assume the transformation and its symmetry group a priori and are strongest exactly where that assumption is right. The complementary situation --- an opaque, already-trained network whose transformation structure is not known --- is the one this paper addresses: we do not build the action, we measure whether one exists, of what form, and at what price, and the measurement then serves as the blueprint for an interface (\S{}7). The two directions meet where a measurement produces a specification: the fixed block-diagonal action of \S{}7.2 is not an architectural prior but a measured premise, and carrier sufficiency (Thm 5) states when such a compiled action is sufficient. There is also a difference in what is claimed: an equivariant architecture asserts a symmetry of the \textit{whole} network, whereas a measurement report can be, and here often is, weaker and more specific --- this action, on this site, for this consumer, with this defect.

\subsection{Colour and perceptual transformations}

Colour and scale have a mature analytic theory of what their transformations are: the axioms of scale-space [17], [18] fix the diffusion family this paper uses as its dissipative instance, and colour-appearance and HSV-family coordinates fix the hue rotation that supplies its group instance. We treat these as the implemented physics of the rendered domain --- a hue rotation and a diffusion semigroup that are exact at the pixel level (Assumption A1) --- and take no position on perceptual uniformity or spectral rendering.

The distinction that matters for us is algebraic rather than perceptual, and it is where this line is silent: the same perceptual space contains a \textit{group} (hue), a \textit{semigroup} (diffusion), and \textit{monoids with degenerate points} (clipped scalings), and the three behave differently under composition. Prior colour-equivariant work targets recognition performance under colour shift; this paper targets the transformation law itself --- its orbit structure, its degeneracies, and what of it remains operable --- and uses the algebra, not the attribute, as the independent variable. That choice is what makes the dissipative family's failure informative rather than embarrassing: the heat semigroup composes exactly at the truth level, and the obstruction appears only when we ask a representation to carry it (\S{}4.7).

\subsection{Causal abstraction and latent linearisation}

The fourth line asks when a high-level intervention corresponds to a low-level one. Causal abstraction [19], [20] gives the interventional-consistency condition; our fibre condition is, formally, that condition specialised to transformation actions. The latent-linearisation line asks when a single direction is a well-defined intervention --- the linear representation hypothesis and its geometry [21], the observation that steering vectors can be steerable but not decodable [22], and the non-identifiability of steering vectors without a stated reference [23], surveyed for representation engineering by Wehner et al. [24] --- with certification-style evaluation proposed for interventional claims [25] and paired-scene benchmarks for compositional faithfulness [26]. Koopman-style approaches [27] ask when a nonlinear system admits an invariant finite-dimensional subspace on which the evolution is linear, and design or learn a lift into it; our band-limited global family is a truncated instance of that lift applied to a \textit{physical} transformation of a frozen vision encoder. What that line does not supply --- and what this paper adds --- is an existence criterion for a \textit{given} representation rather than a designer's lift, together with the transfer from temporal dynamics to non-temporal transformations of images.

Specialising to transformation actions buys three things this line does not contain. First, the \textit{whole algebra} must descend, not single interventions --- and the composition law is then priced quantitatively by the accumulation law, which an intervention-at-a-time condition cannot express. Second, identifiability becomes an object rather than a caveat: the kernel and stabilisers of the representation name exactly which part of the algebra survives, and lumpability of a partition is the classical version of the same condition [28], with the quotient construction itself standard [29], [30]. Third, the question continues past existence to realisation form, family capacity, and construction: we do not assume a linear realisation exists, we measure when it does (\S{}5.2), prove that the natural per-dimension class cannot realise a closed orbit (Thm 3), and build the working object from the measurements (\S{}7).

\subsection{What none of the four lines contains}

The gap is a gap of \textit{question}, not of technique, and it can be said in four sentences. Equivariance measurement answers how much of a transformation a representation carries, but not whether an action exists, which part of the algebra survives, or what a chain of transformations costs. Equivariant architectures answer how to build a prescribed symmetry in, but say nothing about the structure an already-trained opaque network carries. Colour and scale-space theory fix what the reference transformations analytically are, but not whether a learned representation carries them, nor at what price. Causal abstraction and latent linearisation answer when an intervention is well defined, but at the level of single interventions: they contain neither the whole algebra nor its composition price, and they stop before realisation form and construction.

The join is the paper's contribution: a compatibility condition that decides existence, a decomposition that prices realisation, a measured organisation that fixes the form of the operator, and a construction whose action is fixed by that measurement rather than fitted. Appendix G states the nearest collision against each of our claims.
